\documentclass[conference]{IEEEtran}
\IEEEoverridecommandlockouts
% The preceding line is only needed to identify funding in the first footnote. If that is unneeded, please comment it out.
\usepackage{cite}
\usepackage{amsmath,amssymb,amsfonts}
\usepackage{algorithmic}
\usepackage{graphicx}
\usepackage{textcomp}
\usepackage{xcolor}
\usepackage{url}
\def\BibTeX{{\rm B\kern-.05em{\sc i\kern-.025em b}\kern-.08em
    T\kern-.1667em\lower.7ex\hbox{E}\kern-.125emX}}
\begin{document}

\title{Benchmarking TurboQuant and KV Cache Compression Methods on Vietnamese Large Language Models}

\author{
\IEEEauthorblockN{
Do Kien Hung, Phan Trong Qui, Ho Viet Anh, Pham Minh Quan, Tran Minh Khanh, \\
Nguyen Van Quang Duy, Nguyen Ho Phat, Huynh Huu Huy, Huynh Ngoc Thach, \\
Nguyen Dang Quoc Anh, and Phan Trong Phu
}
\IEEEauthorblockA{\textit{Department of Computer Engineering} \\
\textit{Ho Chi Minh City University of Technology and Engineering (HCM-UTE)}\\
Ho Chi Minh City, Vietnam \\
Email: \{group1\}@student.ute.edu.vn}
}

\maketitle

\begin{abstract}
Deploying Large Language Models (LLMs) for real-world Vietnamese applications is growing rapidly, but the Key-Value (KV) Cache scales linearly with context length, leading to severe memory bottlenecks on hardware. While KV Cache compression techniques like TurboQuant, PolarQuant, HQQ, and FP8 show significant promise, their effectiveness has not been systematically evaluated on Vietnamese LLMs where tokenization and linguistic structures differ significantly from English. This paper presents a systematic, reproducible empirical benchmark evaluating the trade-offs between memory footprint, inference latency, throughput, and output quality (perplexity) across various context lengths (from 4k to 32k tokens) for Vietnamese LLMs. Our results establish the Pareto-optimal boundaries for production-level deployments in resource-constrained environments.
\end{abstract}

\begin{IEEEkeywords}
KV Cache Compression, Large Language Models, TurboQuant, PolarQuant, Vietnamese NLP, Hardware Efficiency, Benchmark
\end{IEEEkeywords}

\section{Introduction}
\label{sec:intro}

The deployment of Large Language Models (LLMs) has revolutionized natural language processing (NLP) applications worldwide, including in Vietnam. However, during autoregressive decoding, the Key-Value (KV) Cache grows linearly with both the context length and the batch size, presenting a critical memory bottleneck for serving engines. 

This issue is particularly pronounced for Vietnamese LLMs. Due to tokenization differences and morphological structures, Vietnamese sequences often yield a larger number of tokens per semantic unit compared to English, inflating the KV Cache overhead. This paper establishes a systematic benchmark to evaluate state-of-the-art Post-Training Quantization (PTQ) and compression methods, such as TurboQuant and PolarQuant, on Vietnamese LLMs to determine their hardware efficiency and output quality trade-offs.

\section{Methodology and Experimental Setup}
\label{sec:methodology}

This section describes the theoretical framework of the benchmarked Key-Value (KV) Cache compression algorithms, the hardware serving infrastructure, and the evaluation metrics used to quantify performance and quality trade-offs on Vietnamese LLMs.

\subsection{Theoretical Foundations of KV Cache Compression}

In Transformer-based Large Language Models, the autoregressive generation phase stores the key and value states of all past tokens to avoid recomputation. For a sequence of length $L$, batch size $B$, number of attention heads $H$, and head dimension $D$, the memory footprint of the KV Cache scales as $\mathcal{O}(2 \times B \times L \times H \times D)$. At large context lengths (e.g., $16k$ or $32k$ tokens), the KV Cache dominates the GPU VRAM consumption, leading to severe serving constraints or Out-Of-Memory (OOM) errors.

Post-Training Quantization (PTQ) reduces the memory footprint of the KV Cache by mapping high-precision 16-bit floating-point (BF16/FP16) key-value states to low-bit representations (e.g., 8-bit, 4-bit, or 3-bit). We detail the baseline and advanced quantization methods evaluated in this study below.

\subsubsection{Rotary Position Embedding (RoPE) and Cartesian Outliers}
Modern LLMs, such as LLaMA and Qwen, incorporate Rotary Position Embedding (RoPE) to encode relative positional information. For a 2D coordinate pair $(x_{2j-1}, x_{2j})$ of a key vector $\mathbf{k} \in \mathbb{R}^D$ at position $m$, RoPE applies a rotation matrix:
\begin{equation}
R_{\theta_j, m} = \begin{pmatrix} \cos m\theta_j & -\sin m\theta_j \\ \sin m\theta_j & \cos m\theta_j \end{pmatrix}
\end{equation}
While RoPE enables effective relative position encoding, it introduces a major challenge for standard quantization. The rotation causes key embeddings to exhibit a highly non-uniform distribution with severe coordinate outliers when represented in standard Cartesian coordinates. Applying standard linear quantization in the Cartesian space leads to large quantization errors and subsequent language quality degradation.

\subsubsection{PolarQuant: Polar Coordinate Transformation}
To address the coordinate outliers induced by RoPE, PolarQuant \cite{han2025polarquant} transforms the Cartesian coordinate pairs of the key vectors into polar coordinates. For each pair $(k_{2j-1}, k_{2j})$ of the key vector $\mathbf{k}$, the polar transformation computes the magnitude $r_j$ and the phase angle $\phi_j$:
\begin{equation}
r_j = \sqrt{k_{2j-1}^2 + k_{2j}^2}
\end{equation}
\begin{equation}
\phi_j = \text{atan2}(k_{2j}, k_{2j-1})
\end{equation}
The key insight of PolarQuant is that the phase angles $\phi_j$ are uniformly distributed in the interval $[-\pi, \pi]$ and are highly stable across different sequences and contexts. Consequently, $\phi_j$ can be quantized uniformly using a simple uniform quantizer without storing expensive scaling parameters (scale and zero-point) per block. The magnitude $r_j$ is positive-valued and is quantized using a non-uniform Lloyd-Max scalar quantizer tailored to its empirical distribution.

\subsubsection{TurboQuant: Random Orthogonal Rotation}
TurboQuant \cite{zandieh2025turboquant} expands on the polar coordinate approach by applying a random orthogonal rotation matrix $\mathbf{U} \in \mathbb{R}^{D \times D}$ to the input key vectors prior to quantization:
\begin{equation}
\mathbf{k}' = \mathbf{U} \mathbf{k}
\end{equation}
According to the Central Limit Theorem, applying a random rotation concentrates the coordinate distributions of $\mathbf{k}'$ towards a Gaussian or Beta distribution, eliminating directional outliers and coordinate correlations. This coordinate concentration allows scalar quantization to achieve a near-optimal mean-squared error (MSE) distortion rate without needing complex calibration datasets.

\subsubsection{Quantized Johnson-Lindenstrauss (QJL) Error Correction}
Although random rotations minimize quantization error, low-bit quantization (e.g., 3-bit or 4-bit) still introduces non-negligible error in estimating the query-key inner products, which directly affects the attention logits. To compensate for this estimation error $\epsilon = \mathbf{q}^T \mathbf{k} - \tilde{\mathbf{q}}^T \tilde{\mathbf{k}}$ (where $\tilde{\mathbf{q}}$ and $\tilde{\mathbf{k}}$ are the quantized vectors), TurboQuant integrates a Quantized Johnson-Lindenstrauss (QJL) correction mechanism. 

The residual error vector $\mathbf{v} = \mathbf{k} - \tilde{\mathbf{k}}$ is projected using a random 1-bit projection matrix $\mathbf{S}$:
\begin{equation}
s(\mathbf{v}) = \text{sign}(\mathbf{S} \mathbf{v})
\end{equation}
The query-key attention logit is corrected online during decoding by computing:
\begin{equation}
\text{Attention}(Q, K) \approx \tilde{\mathbf{q}}^T \tilde{\mathbf{k}} + \alpha \cdot \text{sign}(\mathbf{S} \mathbf{q})^T \text{sign}(\mathbf{S} \mathbf{v})
\end{equation}
where $\alpha$ is a scaling constant. The 1-bit representation of the correction vector requires only 1 bit per dimension of VRAM and can be processed using high-speed bitwise operations, correcting attention errors with minimal hardware overhead.

\subsection{Hardware and Serving Infrastructure}
All experiments in this benchmark are executed on an enterprise-grade cloud GPU node to ensure reproducibility:
\begin{itemize}
    \item \textbf{GPU:} NVIDIA RTX 4090 (Ada Lovelace architecture) with 24 GB of GDDR6X VRAM.
    \item \textbf{Operating System:} Ubuntu 22.04 LTS.
    \item \textbf{CUDA Environment:} CUDA Toolkit 12.2 with NVIDIA Driver version 535.104.
    \item \textbf{Serving Engine:} vLLM v0.7.0, utilizing PagedAttention \cite{kwon2023pagedattention} with a logical block size of 16 to manage memory allocations dynamically.
\end{itemize}

\subsection{Benchmarked Models and Dataset}
The benchmark evaluates three high-performance Vietnamese LLM variants:
\begin{enumerate}
    \item \textbf{PhoGPT-7B5-Instruct:} A foundational model trained on 102 billion Vietnamese tokens, using ALiBi positional embeddings \cite{nguyen2023phogpt}.
    \item \textbf{Qwen2.5-7B-Instruct:} A state-of-the-art multilingual model with an native 128k context window and excellent Vietnamese capabilities \cite{qwen2024qwen25}.
    \item \textbf{URA-Llama-3-8B-Instruct:} A Vietnamese instruction-tuned model built on Llama-3, developed by the URA research group at HCMUT.
\end{enumerate}
The test suite consists of `test_set_small.json`, a curated collection of long-context Vietnamese documents spanning $4k$, $8k$, and $16k$ token lengths, processed using NVIDIA NeMo Curator to filter out low-quality text and near-duplicates.

\subsection{Evaluation Metrics}

We evaluate the models across performance (latency and throughput), hardware efficiency (VRAM), and linguistic quality (perplexity).

\subsubsection{Time to First Token (TTFT)}
TTFT measures the time elapsed from submitting the prompt to receiving the first generated token, quantifying prefill phase latency:
\begin{equation}
\text{TTFT} = t_{\text{first\_token}} - t_{\text{prompt\_received}} \quad (\text{ms})
\end{equation}

\subsubsection{Inter-Token Latency (ITL)}
ITL measures the average time taken to generate successive tokens during the decoding phase:
\begin{equation}
\text{ITL} = \frac{1}{M - 1} \sum_{i=2}^{M} (t_i - t_{i-1}) \quad (\text{ms/token})
\end{equation}
where $M$ is the total number of decoded tokens.

\subsubsection{Generation Throughput}
Throughput measures the overall generation speed in tokens per second:
\begin{equation}
\text{Throughput} = \frac{M}{\Delta T} \quad (\text{tokens/s})
\end{equation}
where $\Delta T$ is the total time spent in the decoding phase.

\subsubsection{Perplexity (PPL)}
Perplexity serves as the primary linguistic metric to evaluate text quality degradation under quantization. It is defined as the exponentiated cross-entropy loss over a sequence of $N$ tokens:
\begin{equation}
\text{PPL} = \exp\left(-\frac{1}{N} \sum_{i=1}^{N} \ln P(x_i \mid x_{<i})\right)
\end{equation}
where $P(x_i \mid x_{<i})$ is the probability assigned by the model to the target token $x_i$ given the preceding context.

\subsubsection{GPU Memory Efficiency Index (GMEI)}
To capture the joint efficiency of serving speed and VRAM overhead, we define GMEI as the ratio of decoding throughput to peak dynamic VRAM:
\begin{equation}
\text{GMEI} = \frac{\text{Throughput}}{\text{Peak Dynamic VRAM (MB)}} \quad (\text{tokens/s/MB})
\end{equation}
where Peak Dynamic VRAM is computed by subtracting the static model loading memory from the peak VRAM recorded during generation.

\subsubsection{KV Cache Compression Ratio (KCCR)}
KCCR measures the relative VRAM savings achieved by the compressed format compared to the baseline 16-bit format:
\begin{equation}
\text{KCCR} = \frac{\text{Size}_{\text{compressed\_KV}}}{\text{Size}_{\text{BF16\_KV}}} \times 100\% \quad (\%)
\end{equation}

\subsubsection{Downstream Accuracy Metrics}
For downstream tasks (such as extractive Q\&A from VMLU and SQuAD), we report the F1-Score and Exact Match (EM) metrics. Exact Match is a binary indicator:
\begin{equation}
\text{EM} = \begin{cases} 1 & \text{if } \text{prediction} = \text{ground truth} \\ 0 & \text{otherwise} \end{cases}
\end{equation}
and the F1-score is computed based on token-level precision and recall.

\section{Experimental Evaluation and Results}
\label{sec:results}

In this section, we present the empirical results of our benchmarking study. We compare the performance of various KV Cache compression techniques against the uncompressed 16-bit baseline across different context lengths and models. Detailed performance logs and Pareto trade-off curves are presented to highlight the efficiency gains and quality impact.

\section{Discussion and Future Work}
\label{sec:discussion}

The experimental findings reveal key trade-offs between hardware optimization and linguistic quality when compressing the KV Cache of Vietnamese LLMs. In this section, we analyze the sensitivity of the Vietnamese language to deep cache compression and discuss potential mitigation strategies. Future work will investigate dynamic compression ratios and integration with advanced attention kernels.


\bibliographystyle{IEEEtran}
\bibliography{references}

\end{document}
